% Created 2016-02-28 Sun 15:47
\documentclass[11pt]{article}
\usepackage[utf8]{inputenc}
\usepackage[T1]{fontenc}
\usepackage{fixltx2e}
\usepackage{graphicx}
\usepackage{longtable}
\usepackage{float}
\usepackage{wrapfig}
\usepackage{rotating}
\usepackage[normalem]{ulem}
\usepackage{amsmath}
\usepackage{textcomp}
\usepackage{marvosym}
\usepackage{wasysym}
\usepackage{amssymb}
\usepackage{hyperref}
\tolerance=1000
\author{Giulio Bottazzi}
\date{\today}
\title{Financial Economics 2015/2016}
\hypersetup{
  pdfkeywords={},
  pdfsubject={},
  pdfcreator={Emacs 24.3.1 (Org mode 8.2.4)}}
\begin{document}

\maketitle
\tableofcontents


\section{Structure of the course}
\label{sec-1}

\begin{itemize}
\item Teaching load: 63 hours (ECTS credits 9)
\item Main lecturer: BOTTAZZI Giulio
\item Teaching assistant: GIACHINI Daniele
\item Guest lecturer and experts: TBA
\item Semester: Spring
\item Keywords: financial economics, general equilibrium, arbitrage,
asset pricing, optimization

\item Outline: This course aims to provide a wide theoretical foundation
to the most widespread notions and tools in modern finance,
including present value theory, real interest rates, arbitrage
pricing, portfolio allocation and mean-variance analysis.

\item Objectives: The students should understand the theoretical
justification of interest rate and yield curve. The students
should become acquainted with the notion of arbitrage and
equilibrium prices in different market settings. They should be
able to solve the problem of portfolio optimization and
mean-variance analysis in rather general terms.

\item Description: General equilibrium, no arbitrage in risk less
economies, real interest rate and the interest rate
curve. Arbitrage in security markets, state prices, complete and
incomplete markets, valuation functional and the fundamental
theorem of finance. Choices under uncertainty, expected utility
theory and risk aversion. Portfolio choices, optimal portfolio
with multiple risky assets. General equilibrium under
uncertainty. Pricing kernel and mean-variance analysis. OPTIONAL:
behavioral finance, asset prices under ambiguity, evolutionary
finance and the market selection hypothesis.

\item Prerequisites: The course requires a knowledge of linear algebra
(linear space, linear map, basis, inversion, eigenvectors and
eigensystems), calculus (differential analyses of function of many
real variables and static optimization) and, to a lesser extent,
probability theory. Some basic knowledge of the theory of choice
under uncertainty (expected utility theory) and economic
equilibrium are in general considered understood and only
cursorily reviewed.

\item Final exam: the get the final grading students are required to
pass a final written exam. At least one trial exam will be
organized to test your level of preparation. Students get extra
points for completing their homework during the course.

\item Course web page: \url{http://cafim.sssup.it/~giulio/teaching.html}
\end{itemize}
\section{Course material}
\label{sec-2}

Lecture notes will be made available to all students. The syllabus
(see below) contains references to specific parts of the books. Each
week, students will be assigned a few examples and exercises as
homework. The following week, the teaching assistant will collect
the students papers and review the exercises.

\begin{itemize}
\item Textbooks:
\begin{description}
\item[{[LRW]}] Principles of Financial Economics, S.F. LeRoy and
J. Werner
\item[{[FBJ]}] Foundations of Financial Markets and Institutions,
Fabozzi, Modigliani and Jones
\end{description}

\item Further references and Optional readings:
\begin{description}
\item[{[MWG]}] Microeconomic Theory, A. Mas-Colell, M. D. Whinston,
and J. R. Green
\item[{[CN]}] Quantitative Financial Economics, K. Cuthbertson and
D. Nitzsche
\item[{[HR]}] Financial Economics, T. Hens and M. O. Rieger
\item[{[ES]}] Ambiguity and Asset Markets, L. G. Epstein and
M. Schneider, NBER paper;
\item[{[Hu]}] Options, Futures and Other Derivatives, J. C. Hull (7th ed.)
\item[{[Ni]}] The Ascent of Money: A Financial History of the World,
Niall Ferguson
\item[{[Eco]}] Guide to the Financial Markets, The Economist
\end{description}

\item Prerequisites suggested reading (only relevant parts):
C. P. Simon, L. E. Blume, Mathematics for Economists, H. R. Varian
Microeconomic Analysis.
\end{itemize}

\section{Syllabus}
\label{sec-3}

List of topics with reference to the textbooks. Each item roughly
corresponds to one or two lectures.

\begin{itemize}
\item Introduction: the structure of the course and its textbooks; the
role of financial institutions \href{https://docs.google.com/presentation/d/1fPWtpCZgLpRXPdWiiNIE9HtLyaMs5hSaSQqVL_Gj7No/edit?usp=sharing}{slides} (Textbooks sections: [FBJ]
Ch. 1, 2 and 3; [Ni] Ch. 1; [Eco] Ch 1). The first building
blocks: Walras market clearing condition and utility theory.
\item The notion of economic equilibrium; Two agents, two goods, one
date model. The role of prices.
\item Two date model. Spot and forward market. Financial
equilibrium. Present value prices and economic
equilibrium. Redundant contracts; the role of arbitrage in setting
princes;
\item Real and nominal interest rate; The role of inflation; (Textbooks
sections: [FBJ] Ch. 4 and 5)
\item A second look at intertemporal allocation in a one good two date
model: the role of preferences and endowments. Pricing of future
contracts (Textbooks sections: [Hu] Sections from 4.1 to 4.8 and
from 6.1 to 6.3 [Ni] Ch. 2)
\item The multi-period case and the role of the yield curve. Pricing
loans and fixed income securities.
\item Structure of security markets with uncertainty; complete markets
and redundant assets; (Textbook sections: [LRW] 1.1, 1.2)
\item Structure of security markets: Left and right matrix inverse;
payoff pricing correspondence; the Law of One Price; state claims
and state prices (Textbook sections: [LRW] 1.6, 2.1, 2.3, 2.4,
2.5)
\item Arbitrage and valuation: strong and weak arbitrage; positive
pricing; the payoffs convex cone; Farkas' Lemma; the valuation
functional and strong arbitrage; Stiemke's Lemma; the valuation
functional and arbitrage; (Textbook sections: [LRW] 3.2, 3.3, 3.4,
3.5, 5.1)
\item Arbitrage and valuation: the fundamental theorem of finance;
multiplicity of valuation functional; values bound; risk neutral
probabilities; (Textbook sections: [LRW] 5.2, 5.3, 5.5, 6.3, 6.4,
6.5, 6.6, 6.7;)
\item Valuation of derivatives in a two date model: forward and future;
put and call options. Practical example: from option price to
risk-neutral measure (The software used is \href{http://www.gnu.org/software/octave/}{octave}. Data courtesy
of \href{http://finance.yahoo.com/}{yahoo!finance}) (Textbook sections: see the discussion of
forward and future contracts in [Hu] Ch.1-3)
\item Brief history of risk and uncertainty in Economics
\item Choices under uncertainty: consumption and portfolio choices,
Expected Utility Theory, separable and time invariant utilities;
risk aversion; absolute and relative risk aversion coefficients;
certainty equivalent; concave transformation; comparison of
utility functions (Textbook sections: [LRW] 9.2, 9.3, 9.4, 9.5,
9.6, 9.8; the proves of the theorems can also be found in [MWG]
Ch.6 Sec.B; further examples are discussed in [CN] Ch.1 Sec.3-5)
\item Portfolio optimization: first order conditions; arbitrage and
optimal portfolio, optimal consumption, prices and
returns. (Textbook sections: [LRW] 1.4, 1.5, 2.6, 3.6, 11.1, 11.2;
see [CN] Ch.15 for an introductory discussion of the problem)
\item Portfolio optimization with one risky security: impact of wealth
on optimal investment; exercises (Textbook sections: [LRW ]12.2)
\item Portfolio optimization with several risky securities: from
consumption to investment optimization; risk-return trade off;
optimal portfolio under fair pricing; (Textbook sections: [LRW]
13.2, 13.3, 13.4, 9.9, 13.6)
\item Economic equilibrium with uncertanty; representative agent model;
equilibrium sufficient conditions; equilibrium and arbitrage
(Textbook sections: [LRW] 1.7, 1.8, 1.9, 2.4, 3.7)
\item Economic equilibrium: example with heterogeneous agents,
log-utilities and Arrow securities.
\item Mean variance analysis and the Markowitz economy
\item Consumption based security pricing; a first derivation of CAPM;
the social planner problem; equilibrium and Pareto optimality,
first and second welfare theorems; (Textbook sections: 14.2, 14.3,
14.5, 15.2, 15.3, 15.4, 15.5, without theorem 15.5.1)
\item Practical example: mean-variance efficient frontier from empirical
data; (The software used is \href{http://www.gnu.org/software/octave/}{octave}. Data courtesy of
\href{http://finance.yahoo.com/}{yahoo!finance})
\item CAPM (Textbook sections: chapter 19 all, theorems without proof; ;
see also the discussion in [CN] Ch.15 Sec. 3-5)
\item A short introduction to continuous time option pricing (by a guest
lecturer)
\item Two seminars on empirical problems in finance (guest experts)
\end{itemize}
% Emacs 24.3.1 (Org mode 8.2.4)
\end{document}